% !TeX root = ../../../tesis.tex

% =====================================================================
\subsubsection{El problema de la compensabilidad}
\label{subsubsec:garibelt-compensabilidad}
% =====================================================================

Los cinco indicadores que integra este instrumento no son cinco formas
distintas de medir lo mismo: son cinco dimensiones de un fenómeno que no
admite una unidad de medida única. La cobertura~$C$ opera sobre fracciones
del territorio; el tiempo promedio~$T$ recoge la experiencia del usuario en
minutos; la intermediación~$B(v)$ describe la posición de cada nodo en la
estructura del grafo. Reunir estas magnitudes en un solo número es
técnicamente posible, pero produce una cifra que oscurece más de lo que
revela. La Clasificación Garibelt propone una salida distinta: un tablero de
cinco lecturas donde cada dimensión conserva su propio rango y su propio
significado, y donde las diferencias estructurales permanecen visibles.

Agregar cinco indicadores en un promedio parece resolver el problema de la
heterogeneidad, pero introduce otro: el sesgo de compensabilidad. Una red
puede obtener una nota global «aceptable» si dos dimensiones son excelentes,
aunque una tercera esté en estado crítico; el promedio absorbe ese dato en
lugar de exhibirlo \citep{hossain2024}. \citet{deng2023} documentan este
riesgo en redes multimodales: la concentración de flujo en nodos clave no
queda visible en índices compuestos, y esa invisibilidad dificulta la
priorización de intervenciones. La
Figura~\ref{fig:compensabilidad-garibelt} ilustra el contraste con la
Red~$\beta$: el índice ponderado arroja $0.56$ —Aceptable— mientras el
Espectro MJG deja en primer plano la dimensión de eficiencia en estado
Crítico.

\begin{figure}[htbp]
\centering
% Diagrama: Problema de compensabilidad — Clasificación Garibelt
% Usado en: 5-Analisis-Red-Actual/5-6-Clasificacion-Garibelt/5-6-1-Fundamento.tex
% fig:compensabilidad-garibelt
%
% Muestra cómo un índice compuesto (0.56 — Aceptable) oculta la dimensión
% Crítica que el Espectro Garibelt sí deja visible.
% Colores de bandas: red!65 / orange!65 / yellow!55!green!35 / green!50
% Coordenadas diseñadas para caber en \textwidth (≈14 cm).

\begin{tikzpicture}[
    casilla/.style = {
        rectangle,
        minimum width  = 3.8cm,
        minimum height = 0.65cm,
        inner sep      = 5pt,
        align          = center,
        font           = \small,
        rounded corners= 3pt
    },
    etiq/.style = {
        font       = \small,
        align      = right,
        text width = 3.4cm
    }
]

% ================================================================
%  PANEL IZQUIERDO — índice compuesto
%  Barra de escala [0,1] con 4 bandas; marcador en 0.56
%  Barra: x in [0, 5];  y in [4.10, 4.85]
%  Bandas: 0–1.25 / 1.25–2.50 / 2.50–3.75 / 3.75–5.00
% ================================================================

\node[font=\small\bfseries, align=center] at (2.5, 5.65)
    {Índice compuesto};
\node[font=\scriptsize\itshape, color=gray!60, align=center] at (2.5, 5.25)
    {Red hipotética $\beta$};

% Bandas coloreadas
\fill[red!65]             (0.00, 4.10) rectangle (1.25, 4.85);
\fill[orange!65]          (1.25, 4.10) rectangle (2.50, 4.85);
\fill[yellow!55!green!35] (2.50, 4.10) rectangle (3.75, 4.85);
\fill[green!50]           (3.75, 4.10) rectangle (5.00, 4.85);
\draw[gray!50, thin]      (0.00, 4.10) rectangle (5.00, 4.85);
\draw[gray!55, thin] (1.25, 4.10) -- (1.25, 4.85);
\draw[gray!55, thin] (2.50, 4.10) -- (2.50, 4.85);
\draw[gray!55, thin] (3.75, 4.10) -- (3.75, 4.85);

% Etiquetas de banda
\node[font=\tiny\bfseries, white]  at (0.625, 4.475) {Crít.};
\node[font=\tiny\bfseries]         at (1.875, 4.475) {Débil};
\node[font=\tiny\bfseries]         at (3.125, 4.475) {Acept.};
\node[font=\tiny\bfseries, white]  at (4.375, 4.475) {Idón.};

% Marcador en 0.56 → x = 0.56 × 5 = 2.80
\fill[unamAzul!80]
    (2.80, 4.10) -- (2.67, 3.86) -- (2.93, 3.86) -- cycle;
\node[font=\footnotesize\bfseries, anchor=north]
    at (2.80, 3.85) {$0.56$ — Aceptable};

% Ticks de escala inferior
\foreach \xval/\lab in {0.00/0.0, 1.25/0.25, 2.50/0.50, 3.75/0.75, 5.00/1.0}{
    \draw[gray!45] (\xval, 4.10) -- (\xval, 3.97);
    \node[font=\tiny, color=gray!55, anchor=north] at (\xval, 3.97) {\lab};
}

% ================================================================
%  FLECHA CENTRAL
% ================================================================
\draw[-Stealth, thick, color=gray!45]
    (5.30, 4.48) -- (6.10, 4.48)
    node[midway, above, font=\tiny\itshape, color=gray!55] {desagregado};

% ================================================================
%  PANEL DERECHO — Espectro Garibelt
%  Labels (right-aligned, text width 3.4 cm): centradas en x = 8.0
%  Casillas (width 3.8 cm): centradas en x = 11.80
%  Filas: y = 4.90, 4.08, 3.26, 2.44, 1.62  (dy = 0.82)
% ================================================================

\node[font=\small\bfseries, align=center] at (10.30, 5.65)
    {Espectro Garibelt};
\node[font=\scriptsize\itshape, color=gray!60, align=center] at (10.30, 5.25)
    {misma Red $\beta$};

% Fila 1 — Accesibilidad 0.80 Idóneo
\node[etiq] at (8.0, 4.90) {Accesibilidad ($C$)};
\node[casilla, fill=green!50, draw=gray!40]
    at (11.80, 4.90) {$0.80$ — \textbf{Idóneo}};

% Fila 2 — Conectividad 0.75 Idóneo
\node[etiq] at (8.0, 4.08) {Conectividad ($C_i$)};
\node[casilla, fill=green!50, draw=gray!40]
    at (11.80, 4.08) {$0.75$ — \textbf{Idóneo}};

% Fila 3 — Eficiencia 0.20 Crítico (resaltada)
\node[etiq, font=\small\bfseries] at (8.0, 3.26) {Eficiencia ($DI$)};
\node[casilla, fill=red!65, text=white, draw=red!70, line width=0.8pt]
    at (11.80, 3.26) {$0.20$ — \textbf{Crítico}};

% Fila 4 — Fluidez 0.45 Débil
\node[etiq] at (8.0, 2.44) {Fluidez ($T$)};
\node[casilla, fill=orange!65, draw=gray!40]
    at (11.80, 2.44) {$0.45$ — Débil};

% Fila 5 — Centralidad 0.60 Aceptable
\node[etiq] at (8.0, 1.62) {Centralidad ($B(v)$)};
\node[casilla, fill=yellow!55!green!35, draw=gray!40]
    at (11.80, 1.62) {$0.60$ — Aceptable};

% Nota al pie del panel derecho
\node[font=\scriptsize\itshape, color=red!60,
      anchor=north, text width=4.6cm, align=center]
    at (10.30, 1.15)
    {Dimensión Crítica: visible en el espectro,\\
     oculta en el índice compuesto};

\end{tikzpicture}

\caption[Problema de compensabilidad en índices compuestos de transporte]{%
    Problema de compensabilidad en una red hipotética ($\beta$). A la
    izquierda, el índice ponderado produce $0.56$ (Aceptable), ocultando
    que la eficiencia de ruta está en estado Crítico. A la derecha, el
    Espectro MJG preserva esa información al mantener cada dimensión
    en su propia banda.}
\label{fig:compensabilidad-garibelt}
\fuentefigura{Elaboración propia. Ejemplo hipotético con fines ilustrativos.}
\end{figure}

% =====================================================================
\subsubsection{Normalización como calibración, no como agregación}
\label{subsubsec:garibelt-normalizacion}
% =====================================================================

Resolver el problema de las unidades heterogéneas no es lo mismo que
sumarlas: es calibrarlas. Llevar magnitudes distintas a una escala común
sirve para que sean comparables en un mismo tablero, no para operarlas
aritméticamente. \citet{kujala2018} construyeron un repositorio de
indicadores de redes de transporte para 25 ciudades normalizando los datos
\gtfs{} de cada una, pero sin producir un índice global: el resultado de
su trabajo es un perfil multidimensional por ciudad, precisamente porque
normalizar y agregar son operaciones con propósitos distintos. La
Clasificación Garibelt sigue esa misma lógica: normaliza para que las
cinco dimensiones se lean en la misma escala de bandas, no para sumarlas.

% =====================================================================
\subsubsection{Calibraciones distintas para naturalezas distintas}
\label{subsubsec:garibelt-calibraciones}
% =====================================================================

Si cada dimensión tiene propiedades estadísticas distintas, no tendría
sentido normalizarlas con el mismo procedimiento. \citet{talukder2017}
formalizan este criterio al desarrollar el índice de Objetivos de
Desarrollo Sostenible: el Paso~5 de las guías de la \textsc{ocde} exige
que el método de normalización se ajuste tanto al marco teórico de la
variable como a las propiedades de su distribución. En la Clasificación
Garibelt esto se traduce así: $T$ se calibra contra el rango de tiempos
de traslado documentado por la \textsc{enoe} para la \cdmx{} (85–180
min), porque existe una referencia externa confiable; $DI$ se normaliza
contra sus límites teóricos ($1.0$–$2.5$), porque la teoría de grafos
define con precisión cuándo una ruta es óptima y cuándo es
inaceptablemente desviada; $C_i$ usa percentiles del propio dataset,
porque su significado es relativo al comportamiento de la red analizada;
y $B(v)$ se transforma mediante el índice de Gini, que captura si el
flujo de intermediación está concentrado o distribuido.

% =====================================================================
\subsubsection{Por qué estas cinco dimensiones}
\label{subsubsec:garibelt-dimensiones}
% =====================================================================

Detrás de la selección de los cinco indicadores hay un marco teórico
explícito. \citet{bocarejo2012} identifican cuatro componentes de la
accesibilidad urbana que no tienen una unidad de medida en común:
infraestructura, uso de suelo, restricciones temporales y condición del
usuario. Los tres primeros son los que el instrumento captura. $T$ y $DI$
miden qué tan eficiente y rápida es la red para conectar orígenes con
destinos; $B(v)$ señala si esa eficiencia depende de pocos nodos o está
distribuida; $C$ informa qué fracción del territorio metropolitano queda
dentro del radio de caminata de alguna parada; $C_i$ refleja cuántas rutas
convergen en cada nodo, funcionando como proxy de la oferta temporal. La
condición socioeconómica del usuario ---el cuarto componente de Bocarejo y
Oviedo--- queda diferida a la versión~2 del instrumento.
